% Brüche und Dezimalzahlen · Anteil bestimmen · Prüfungs-Fokus MSA – bank/brueche-dezimalzahlen/e1.jsonl (zusammenbau v0.6, Prüfungs-Fokus)
\einheitenkopf[e1]{Einheit 1 · Bruch als Anteil}
\zweigzeile{ab Kl. 5 · P10 ×9}
\setcounter{aufgabe}{0}

% Anlauf (ohne Prüfkennung)
\begin{aufgabe}{Anteil bestimmen – Anlauf}
\begin{teile}
\teil Sind alle Teile gleich groß? \janein

\bruchkreis{0}{6}
\teil Welcher Anteil des Streifens ist gefärbt?

\bruchrechteck{3}{5}
\leerfeld
\teil Färbe $\frac{5}{6}$ des Streifens.

\bruchrechteck{0}{6}
\end{teile}
\end{aufgabe}

\begin{aufgabe}{Anteil bestimmen – Prüfungsaufgaben}
\begin{teile}
\teil Markiert sind die Sektoren A und E. Welcher Anteil des Kreises ist markiert? \hfill (P19)\\ \kreuz{$\frac{2}{5}$}\\ \kreuz{$\frac{3}{9}$}\\ \kreuz{$\frac{2}{9}$}\\ \kreuz{$\frac{3}{5}$}

\kreisdiagramm{A/80,B/80,C/80,D/80,E/40}
\teil Markiert sind die Sektoren A und B. Welcher Anteil des Kreises ist markiert? \hfill (P19)\\ \kreuz{$\frac{2}{6}$}\\ \kreuz{$\frac{3}{8}$}\\ \kreuz{$\frac{2}{8}$}\\ \kreuz{$\frac{3}{6}$}

\kreisdiagramm{A/90,B/45,C/90,D/45,E/45,F/45}
\teil Das Rechteck besteht aus $12$ Kästchen. Die beiden schrägen Linien bilden mit der unteren Seite ein Dreieck. Welcher Anteil des Rechtecks ist das Dreieck? \hfill (P24) \\

\begin{ksys}[xmin=0,xmax=6,ymin=0,ymax=2,ablesen] \funktionab{\x}{}{0}{2} \funktionab{4-\x}{}{2}{4} \end{ksys}
\leerfeld
\teil Das Quadrat besteht aus $16$ Kästchen. Die beiden schrägen Linien bilden mit der oberen Seite ein Dreieck. Welcher Anteil des Quadrats ist das Dreieck? \hfill (P24) \\

\begin{ksys}[xmin=0,xmax=4,ymin=0,ymax=4,ablesen] \funktionab{4-\x}{}{0}{2} \funktionab{\x}{}{2}{4} \end{ksys}
\leerfeld
\teil Welcher Anteil der Fläche ist grau? Gib ihn als Bruch und in Prozent an. \hfill (P14) \\

\bruchrechteck{13}{25}
Bruch: \leerfeld Prozent: \leerfeld[\%]
\teil Welcher Anteil der Figur ist grau? Gib ihn als Bruch und in Prozent an. \hfill (P14) \\

\bruchrechteck{11}{20}
Bruch: \leerfeld Prozent: \leerfeld[\%]
\teil Markiere $30\,\%$ der Fläche. \hfill (P22)

\bruchrechteck{0}{20}
\teil Markiere $40\,\%$ der Fläche. \hfill (P22)

\bruchrechteck{0}{25}
\teil Markiere $25\,\%$ der Figur. \hfill (P23)

\bruchrechteck{0}{32}
\teil Färbe $25\,\%$ der Figur. \hfill (P23)

\bruchrechteck{0}{36}
\teil Schraffiere $\frac{4}{5}$ der Fläche. \hfill (P17)

\bruchrechteck{0}{20}
\teil Schraffiere $\frac{2}{3}$ der Fläche. \hfill (P17)

\bruchrechteck{0}{18}
\teil Markiere $\frac{5}{12}$ des Rechtecks. \hfill (P21)

\bruchrechteck{0}{12}
\teil Markiere $\frac{4}{9}$ des Rechtecks. \hfill (P21)

\bruchrechteck{0}{9}
\teil Kennzeichne ein Achtel der Fläche des Quadrats. \hfill (P25)

\raute{4}{4}
\teil Kennzeichne drei Viertel der Kreisfläche. \hfill (P25)

\kreisleer[2]
\teil Bei welcher Figur ist genau $\frac{1}{6}$ grau? \hfill (P26F)\\ \kreuz{P}\\ \kreuz{Q}\\ \kreuz{R}\\ \kreuz{S}

P \streifen[0]{20}{}{} \quad Q \kreissektor{60}{} \quad R \bruchrechteck{1}{4} \quad S \bruchkreis{1}{3}
\teil Bei welcher Figur ist genau $\frac{1}{5}$ grau? \hfill (P26F)\\ \kreuz{P}\\ \kreuz{Q}\\ \kreuz{R}\\ \kreuz{S}

P \bruchkreis{1}{4} \quad Q \streifen[0]{30}{}{} \quad R \bruchrechteck{1}{6} \quad S \kreissektor{72}{}
\end{teile}
\end{aufgabe}

\medskip
\uebersichtskasten{Bruch: Der Nenner sagt, in wie viele gleich große Teile das Ganze geteilt ist; der Zähler sagt, wie viele Teile gemeint sind. \\ Figur aus 10 gleichen Feldern, 7 gefärbt: Anteil 7/10 \quad 4 von 20 Kindern: 4/20 = 1/5 \\ Bruchteil einer Menge: Ganzes durch den Nenner teilen, dann mal Zähler. \\ 2/5 von 35 €: 35 : 5 = 7 → 2 · 7 = 14 € \\ Rest zum Ganzen: Alle Teile zusammen sind 1. \quad 9/10 voll → 1/10 fehlt}
